\subsection{Morfismos cuasicompactos.}
\label{subsection:IV.1.1-fr}

\begin{definition}[1.1.1]
\label{IV.1.1.1-fr}
Se dice que un morfismo de preesquemas $f:X\to Y$ es cuasicompacto si
la aplicación continua $f$ del espacio topológico $X$ al espacio
topológico $Y$ es cuasicompacta \textup{(\textbf{0}\textsubscript{III}, 9.1.1)}, es
decir, si la imagen inversa $f^{-1}(U)$ de todo abierto cuasicompacto $U$ de $Y$
es cuasicompacta \textup{(cf. (\textbf{I}, 6.6.1))}.
\end{definition}

Si $\mathfrak B$ es una base de la topología de $Y$ formada por abiertos
afines, para que $f$ sea cuasicompacto, es necesario y suficiente que para todo
$V\in\mathfrak B$, $f^{-1}(V)$ sea \emph{unión finita de abiertos afines}.
Por ejemplo, si $Y$ es afín y $X$ cuasicompacto, \emph{todo} morfismo
$f:X\to Y$ es cuasicompacto (\textbf{I}, 6.6.1).

Si $f:X\to Y$ es un morfismo cuasicompacto, está claro que para todo
abierto $V$ de $Y$, la restricción de $f$ a $f^{-1}(V)$ es un morfismo
cuasicompacto $f^{-1}(V)\to V$. Recíprocamente, si $(U_\alpha)$ es un
recubrimiento abierto de $Y$ y $f:X\to Y$ un morfismo tal que las restricciones
$f^{-1}(U_\alpha)\to U_\alpha$ sean cuasicompactas, entonces $f$ es
cuasicompacto. Por
\oldpage[IV-1]{225}
consiguiente, si $f:X\to Y$ es un $S$-morfismo de $S$-preesquemas, y si existe
un recubrimiento abierto $(S_\lambda)$ de $S$ tal que las restricciones
$g^{-1}(S_\lambda)\to h^{-1}(S_\lambda)$ de $f$ (donde $g$ y $h$ son los
morfismos estructurales) sean cuasicompactas, entonces $f$ es cuasicompacto.

\begin{proposition}[1.1.2]
\label{IV.1.1.2-fr}
\begin{enumerate}
  \item[\textup{(i)}] Una inmersión $X\to Y$ es cuasicompacta si es cerrada, o si
  el espacio subyacente a $Y$ es localmente noetheriano o si el espacio
  subyacente a $X$ es noetheriano.
  \item[\textup{(ii)}] El compuesto de dos morfismos cuasicompactos es cuasicompacto.
  \item[\textup{(iii)}] Si $f:X\to Y$ es un $S$-morfismo cuasicompacto, lo mismo ocurre con
  $f_{(S')}:X_{(S')}\to Y_{(S')}$ para toda extensión $g:S'\to S$ del
  preesquema de base.
  \item[\textup{(iv)}] Si $f:X\to X'$ y $g:Y\to Y'$ son dos $S$-morfismos
  cuasicompactos,
  \[
    f\times_S g:X\times_S Y\to X'\times_S Y'
  \]
  es cuasicompacto.
  \item[\textup{(v)}] Si el compuesto $g\circ f$ de dos morfismos $f:X\to Y$, $g:Y\to Z$
  es cuasicompacto, y si $g$ es separado, o el espacio subyacente a $X$
  localmente noetheriano, $f$ es cuasicompacto.
  \item[\textup{(vi)}] Para que un morfismo $f$ sea cuasicompacto, es necesario y suficiente que
  $f_{\mathrm{red}}$ lo sea.
\end{enumerate}
\end{proposition}

Para la demostración, véase (\textbf{I}, 6.6.4). Se observará que la afirmación (vi) es
también consecuencia de la proposición más general:

\begin{proposition}[1.1.3]
\label{IV.1.1.3-fr}
Sean $f:X\to Y$, $g:Y\to Z$ dos morfismos. Si $g\circ f$ es
cuasicompacto y $f$ sobreyectivo, $g$ es cuasicompacto.
\end{proposition}

\begin{proof}
En efecto, si $V$ es un abierto cuasicompacto en $Z$,
$f^{-1}(g^{-1}(V))$ es cuasicompacto por hipótesis, y se tiene
$g^{-1}(V)=f(f^{-1}(g^{-1}(V)))$, puesto que $f$ es sobreyectivo; por tanto
$g^{-1}(V)$ es cuasicompacto.
\end{proof}

\begin{corollary}[1.1.4]
\label{IV.1.1.4-fr}
Sean $f:X\to Y$, $g:Y'\to Y$ dos morfismos, $X'=X\times_Y Y'$,
$f'=f_{(Y')}:X'\to Y'$; supongamos que $g$ es cuasicompacto y sobreyectivo \textup{(resp.}
sobreyectivo\textup{)}. Entonces, para que $f$ sea cuasicompacto \textup{(resp.} dominante\textup{)}, basta
que $f'$ lo sea.
\end{corollary}

\begin{proof}
Si $g':X'\to X$ es la proyección canónica, $g'$ es un morfismo sobreyectivo
(\textbf{I}, 3.5.2), y se tiene $f\circ g'=g\circ f'$; si $f'$ es cuasicompacto
(resp. dominante), lo mismo ocurre con $g\circ f'$ puesto que $g$ es
cuasicompacto (1.1.2) (resp. sobreyectivo); como $f\circ g'$ es cuasicompacto
(resp. dominante) y $g'$ es sobreyectivo, se deduce que $f$ es
cuasicompacto (1.1.3) (resp. dominante).
\end{proof}

Para abreviar, llamaremos \emph{punto maximal} de un preesquema $X$ a todo
punto genérico de una componente irreducible de $X$; si
$X=\operatorname{Spec}(A)$ es afín, los puntos maximales de $X$ son los
\emph{ideales primos minimales} de $A$.

\begin{proposition}[1.1.5]
\label{IV.1.1.5-fr}
Sea $f:X\to Y$ un morfismo cuasicompacto. Las condiciones siguientes son
equivalentes:
\begin{enumerate}
  \item[\textup{a)}] $f$ es dominante;
  \item[\textup{b)}] para todo punto maximal $y$ de $Y$, se tiene
  $f^{-1}(y)\neq 0$;
  \item[\textup{c)}] para todo punto maximal $y$ de $Y$, $f^{-1}(y)$ contiene un punto
  maximal de $X$.
\end{enumerate}
\end{proposition}

La equivalencia de \emph{a)} y \emph{c)} se demostró en (\textbf{I}, 6.6.5). Está
claro que \emph{c)} implica \emph{b)}; por otra parte, si $X'$ es una
componente irreducible de $X$ que corta a $f^{-1}(y)$, $f(X')$ es
irreducible y contiene a $y$, por lo que está contenido en la componente irreduc%
\oldpage[IV-1]{226}
ible
$Y'=\overline{\{y\}}$ de $Y$ (\textbf{0}\textsubscript{I}, 2.1.6); si $x$ es el
punto genérico de $X'$, se tiene, por tanto, necesariamente $f(x)=y$
(\textbf{0}\textsubscript{I}, 2.1.5); por tanto \emph{b)} implica \emph{c)}.

\begin{proposition}[1.1.6]
\label{IV.1.1.6-fr}
Sean $f':X'\to Y$, $f'':X''\to Y$ dos morfismos, $X$ el preesquema suma
$X'\amalg X''$, $f:X\to Y$ el morfismo igual a $f'$ en $X'$ y a $f''$
en $X''$. Para que $f$ sea cuasicompacto, es necesario y suficiente que $f'$ y
$f''$ lo sean.
\end{proposition}

Esto resulta inmediatamente de la definición.

